% Chapter 5, Topic _Linear Algebra_ Jim Hefferon
%  https://hefferon.net/linearalgebra
%  2021-Jul-28
\topic{Hasil Kali Dalam}
\index{hasil kali dalam|(}
% Smash the depth, not the height; for use in radicals
\def\smashdp#1{{\setbox0=\hbox{$#1$}\dp0=0pt \box0\relax}}

Bab pertama membahas operasi hasil kali titik.
Operasi ini menerima sepasang vektor dari suatu $\R^n$ dan menghasilkan sebuah skalar.
\begin{equation*}
  \colvec{-1 \\ 2 \\ 3}\dotprod\colvec{1 \\ 0 \\ 3}
  =-1\cdot 1+2\cdot 0+3\cdot 3
  =8
\end{equation*}
Dalam bab itu kita belum mengenal istilah `linear', tetapi
operasi ini mempertahankan penjumlahan
\begin{equation*}
 (\vec{v}_1+\vec{v}_2)\dotprod\vec{w}=\vec{v}_1\dotprod\vec{w}
                                      +\vec{v}_2\dotprod\vec{w}
\qquad
 \vec{v}\dotprod(\vec{w}_1+\vec{w}_2)=\vec{v}\dotprod\vec{w}_1
                                      +\vec{v}\dotprod\vec{w}_2
\end{equation*}
dan perkalian skalar.
\begin{equation*}
 (r\vec{v})\dotprod\vec{w}=r(\vec{v}\dotprod\vec{w})
  \qquad
 \vec{v}\dotprod(r\vec{w})=r(\vec{v}\dotprod\vec{w})
\end{equation*}

Satu hal yang menarik tentang operasi hasil kali titik
adalah bahwa operasi ini menggambarkan geometri
$\Re^n\!$.
Sebagai contoh, dua vektor saling tegak lurus jika dan hanya jika
hasil kali titiknya nol, dan panjang sebuah vektor ditentukan oleh
$\vec{v}\dotprod\vec{v}=\absval{\vec{v}\,}^2$.
Di sini kita akan memperluas operasi ini, beserta konsep panjang,
ke ruang-ruang vektor lain.
Kita akan memakai notasi matematika tradisional untuk perluasan ini,
yang tidak berupa titik, melainkan
$\innerprod{\vec{v}}{\vec{w}}$.

Untuk memotivasi definisinya,
kita memerlukan pembahasan bilangan kompleks yang lebih lengkap daripada
yang telah kita tinjau pada bagian pembuka bab ini.
Di sana kita menggambarkan~$\C$ sebagai
himpunan semua $z=a+bi$, dengan $a$ dan~$b$ bilangan real
dan $i^2=-1$.
Kita juga membahas penjumlahan, perkalian skalar,
pengurangan, perkalian, serta panjang.
Untuk pembagian, perhatikan $z=(1+2i)/(3+4i)$.
Rumus selisih kuadrat memberikan
$(a+bi)(a-bi)=a^2-(bi)^2=a^2-(-b^2)=a^2+b^2$,
yang merupakan bilangan real.
Jadi, dengan mengalikan pembilang dan penyebut $z$ dengan $3-4i$,
kita memperoleh bentuk $a+bi$.
\begin{equation*}
  z=\frac{1+2i}{3+4i}\cdot\frac{3-4i}{3-4i}
  =
  \frac{(1\cdot 3+2\cdot 4)+(1\cdot (-4)+2\cdot 3)i}{3^2+4^2}
  =
  \frac{11+2i}{25}
  =
  \frac{11}{25}+\frac{2}{25}\,i
\end{equation*}
Dengan memperhatikan perhitungan tersebut, kita
mendefinisikan \definend{konjugat} bilangan kompleks~$z$
sebagai $\compconj{z}=\compconj{a+bi}=a-bi$, sehingga $z\compconj{z}=\absval{z}^2$.
Membagi persamaan itu dengan $\absval{z}^2$ menghasilkan $z\compconj{z}/\absval{z}^2=1$,
sehingga setiap bilangan kompleks tak nol mempunyai invers perkalian,
$z^{-1}\!=\compconj{z}/\absval{z}^2$.
% (and thus $\absval{z}=1$ if and only if $z^{-1}=\compconj{z}$).
Perhatikan bahwa suatu bilangan kompleks merupakan bilangan real jika dan hanya jika
bilangan itu sama dengan konjugatnya,
bahwa $\compconj{z_1+z_2}=\compconj{z_1}+\compconj{z_2}$,
dan bahwa $\compconj{z_1\cdot z_2}=\compconj{z_1}\cdot\compconj{z_2}$.
Penggambaran $z=a+bi$ pada
bidang kompleks\index{bidang kompleks}
memperlihatkan konjugat tersebut.
\begin{center}
  \includegraphics{jc/asy/innerproduct000.pdf}
\end{center}
Gambar ini memperjelas bahwa
$z+\compconj{z}=(a+bi)+(a-bi)$ sama dengan
$2\cdot a$, yang kita tulis sebagai $2\cdot\RE{z}$.
(``$\RE{z}$'' dibaca sebagai ``bagian real dari $z$.'')
% It also shows the \definend{argument}, $\arg(z)$, the angle that
% $z$ makes with the real axis.

Sekarang kita siap memperumum operasi hasil kali titik dan panjang.
Hasil kali titik dan panjang sudah akrab dalam ruang yang anggotanya berupa
vektor kolom real, tetapi kita telah melihat banyak contoh ruang vektor
yang sangat berbeda, seperti ruang polinomial $\polyspace_n$.
Dalam ruang semacam itu, definisi panjang yang berguna atau alami
tidaklah langsung tampak.
Untuk memperoleh gambaran tentang apa yang harus dilakukan,
perhatikan ruang vektor kompleks~$\C^n$.

Dengan mengingat hasil kali titik pada ruang real,
dugaan pertama yang alami ialah mencoba menjadikan
$\innerprod{\vec{v}}{\vec{w}}$
sebagai kombinasi linear komponen-komponennya, sehingga dalam kasus~$\C^2$
\begin{equation*}
  \vec{v}=\colvec{a_{1,1}+b_{1,1}i \\ a_{2,1}+b_{2,1}i}
  \qquad
  \vec{w}=\colvec{a_{1,2}+b_{1,2}i \\ a_{2,2}+b_{2,2}i}
\end{equation*}
dugaan untuk $\innerprod{\vec{v}}{\vec{w}}$ adalah
$(a_{1,1}+b_{1,1}i)\cdot(a_{1,2}+b_{1,2}i)+(a_{2,1}+b_{2,1}i)\cdot(a_{2,2}+b_{2,2}i)$.
Namun, ini tidak dapat dipakai.
Kita sedang menyelidiki panjang,
jadi kita menginginkan $\innerprod{\vec{v}}{\vec{v}}=\absval{\vec{v}\,}^2\!$,
dan karena itu kita tidak menginginkan $\innerprod{\vec{v}}{\vec{v}}$ sama dengan
$(a_1+b_1i)(a_1+b_1i)+(a_2+b_2i)(a_2+b_2i)$.
Sebaliknya, kita menginginkan sesuatu seperti
$(a_1+b_1i)(a_1-b_1i)+(a_2+b_2i)(a_2-b_2i)$.
Karena itu, kita menggunakan definisi berikut untuk $\C^2\!$.
\begin{equation*}
  \innerprod{\vec{v}}{\vec{w}}
  =(a_{1,1}+b_{1,1}i)\cdot(a_{1,2}-b_{1,2}i)\,+\,(a_{2,1}+b_{2,1}i)\cdot(a_{2,2}-b_{2,2}i)
\end{equation*}
\definend{hasil kali dalam Hermitian}
adalah sebagai berikut:\index{hasil kali dalam!Hermitian}\index{hasil kali dalam Hermitian}
jika vektor $\vec{v},\vec{w}\in\C^n$
mempunyai komponen $v_1$, \ldots~$v_n$ dan $w_1$, \ldots~$w_n$, maka
$\innerprod{\vec{v}}{\vec{w}}=v_1\compconj{w}_1+\cdots+v_n\compconj{w}_n$.
Perhatikan bahwa definisi ini memperluas hasil kali titik dalam arti bahwa,
jika kita kembali ke bilangan real dengan mengambil semua $b$ sama dengan nol,
hasilnya sama dengan $\vec{v}\dotprod\vec{w}$.

Operasi ini menerima dua vektor tetapi menghasilkan sebuah skalar.
Berikut sebuah contoh.
\begin{align*}
  \innerprod{\colvec{1+2i \\ 3+4i}}{\colvec{5+6i \\ 7+8i}}
  &=(1+2i)(5-6i)+(3+4i)(7-8i)                     \\[-2ex]
  &=(17+4i)+(53+4i)=70+8i
\end{align*}

Hasil kali dalam Hermitian bersifat linear pada masukan pertama:~mudah untuk memeriksa bahwa
$\innerprod{\vec{v}_1+\vec{v}_2}{\vec{w}}
 = \innerprod{\vec{v}_1}{\vec{w}}
   +\innerprod{\vec{v}_2}{\vec{w}}$
 dan bahwa
$\innerprod{z\cdot\vec{v}}{\vec{w}}
     =z\cdot\innerprod{\vec{v}}{\vec{w}}$.
Namun, operasi ini tidak linear pada masukan kedua karena, meskipun
$\innerprod{\vec{v}}{\vec{w}_1+\vec{w}_2}
 = \innerprod{\vec{v}}{\vec{w}_1}
   +\innerprod{\vec{v}}{\vec{w}_2}$
berlaku, pada perkalian skalar muncul suatu konjugasi:
$\innerprod{\vec{v}}{z\cdot\vec{w}}
     =\compconj{z}\cdot\innerprod{\vec{v}}{\vec{w}}$.
Tentu saja, perbedaan ini muncul karena dalam definisi tersebut
masukan kedua dikonjugasi, sedangkan masukan pertama tidak.
Akibat lain dari asimetri ini adalah
$\innerprod{\vec{w}}{\vec{v}}\neq \innerprod{\vec{v}}{\vec{w}}$; sebaliknya,
$\innerprod{\vec{w}}{\vec{v}}=\compconj{\innerprod{\vec{v}}{\vec{w}}}$.

Hasil kali dalam Hermitian merupakan salah satu cara, pada kelas ruang
vektor tertentu, untuk memperluas hasil kali titik.
Ada banyak perluasan lain pada banyak ruang lain.
Namun, kita menginginkan perluasan yang menghasilkan gagasan panjang yang wajar.
Untuk itu, \definend{hasil kali dalam} adalah operasi
yang menerima dua vektor dari suatu ruang vektor real atau kompleks,
menghasilkan sebuah skalar, dan memenuhi tiga syarat.
Syarat~(1) menyatakan bahwa operasi itu linear pada masukan pertama, sehingga
$\innerprod{\vec{v}_1+\vec{v}_2}{\vec{w}}
 = \innerprod{\vec{v}_1}{\vec{w}}
   +\innerprod{\vec{v}_2}{\vec{w}}$
dan $\innerprod{z\vec{v}}{\vec{w}}
     =z\innerprod{\vec{v}}{\vec{w}}$.
% (Some presentations take inner product 
% to be linear in the second input; more on this
% at the end of this Topic.)
Syarat~(2) menyatakan bahwa
$\innerprod{\vec{w}}{\vec{v}}=\compconj{\innerprod{\vec{v}}{\vec{w}}}$.
Sebagai akibatnya, perhatikan bahwa
$\innerprod{\vec{v}}{\vec{v}}$ sama dengan konjugatnya dan karena itu
merupakan bilangan real.
Syarat~(3) menyatakan bahwa bilangan real tersebut memenuhi
$\innerprod{\vec{v}}{\vec{v}}\geq 0$ untuk setiap $\vec{v}$, dengan
$\innerprod{\vec{v}}{\vec{v}}= 0$ jika dan hanya jika $\vec{v}=\zero$.

Ruang vektor yang dilengkapi dengan operasi hasil kali dalam disebut
\definend{ruang hasil kali dalam}.
(Pada paragraf sebelumnya kita memakai notasi $\innerprod{\vec{v}}{\vec{w}}$
untuk sebarang hasil kali dalam, sedangkan sebelumnya kita memakainya untuk
kasus khusus hasil kali dalam Hermitian.
Pemakaian ganda ini tidak akan menimbulkan kebingungan karena dalam suatu ruang
hasil kali dalam hanya ada tepat satu operasi semacam itu.)

Setiap ruang vektor real $\R^n$ yang dilengkapi hasil kali titik merupakan
ruang hasil kali dalam.
Hasil kali dalam Hermitian menjadikan setiap ruang vektor kompleks~$\C^n$
sebagai ruang hasil kali dalam.
Pembuktian untuk keduanya terdapat dalam latihan.

Contoh yang tampak sangat berbeda adalah ruang vektor real yang
anggota-anggotanya berupa fungsi bernilai real yang kontinu pada
interval tertutup $\closedinterval{0}{1}$.
Definisikan hasil kali dalamnya sebagai berikut.
\begin{equation*}
  \innerprod{f}{g}=\int_{0}^1 f(t)\cdot g(t)\,dt
\end{equation*}
Operasi ini memenuhi syarat~(1) dalam definisi hasil kali dalam karena
$\int_{0}^1 (f_1(t)+f_2(t))\cdot g(t)\,dt$ sama dengan
$\int_{0}^1 f_1(t)g(t)\,dt+\int_{0}^1 f_2(t)g(t)\,dt$
menurut suatu hasil dari Kalkulus dasar, dan hasil serupa berlaku untuk
perkalian skalar.
Operasi ini memenuhi~(2) karena konjugasi tidak mengubah bilangan real dan
$f(t)g(t)$ sama dengan $g(t)f(t)$.
Untuk syarat~(3), jika~$f$ bukan fungsi nol, maka kuadrat
$f(t)\cdot f(t)$ lebih besar daripada~$0$ untuk suatu
$\hat{t}\in \closedinterval{0}{1}$.
Karena kontinuitas, kuadrat tersebut juga lebih besar daripada nol pada suatu
subinterval di sekitar~$\hat{t}$, sehingga luas di bawah kurvanya lebih besar
daripada~$0$.

Selanjutnya kita memperumum operasi panjang.
Kita akan menganggap suatu operasi sebagai perluasan panjang yang wajar
jika operasi itu menerima vektor, menghasilkan bilangan real, dan memenuhi
tiga syarat; setiap operasi semacam itu kita sebut
\definend{norma},\index{norma} dan kita nyatakan dengan~$\norm{\vec{v}}$.
Syarat pertama ialah $\norm{\vec{v}}\geq 0$ dan
$\norm{\vec{v}}= 0$ jika dan hanya jika $\vec{v}=\zero$.
Syarat kedua ialah bahwa penskalaan ulang suatu vektor juga menskalakan ulang
normanya, secara khusus
$\norm{z\cdot \vec{v}}=\absval{z}\cdot\norm{\vec{v}}$.
Syarat ketiga ialah
Ketaksamaan Segitiga,\index{Ketaksamaan Segitiga}
yaitu $\norm{\vec{v}+\vec{w}}\leq\norm{\vec{v}}+\norm{\vec{w}}$.

Contoh norma yang sudah dikenal dalam ruang vektor~$\R^n$ ialah panjang
Euklides biasa, $\norm{\vec{v}}=\sqrt{\smashdp{v_1^2+\cdots+v_n^2}}$.

Contoh lain terdapat pada ruang vektor real matriks
$\matspace_{\nbym{n}{m}}$.
Definisikan norma matriks, $\norm{M}$, sebagai nilai mutlak terbesar di antara
semua entrinya.
Syarat~(1) jelas, begitu pula syarat~(2).
Syarat~(3) merupakan Ketaksamaan Segitiga biasa untuk bilangan real,
$\absval{a+b}\leq\absval{a}+\absval{b}$.

% Still another example of a norm operation 
% is on the vector space of continuous real-valued functions 
% on the closed interval $\closedinterval{0}{1}$.  
% \begin{equation*}
%   \norm{f}
%   =\int_{x=0}^1 \absval{f(x)}\,dx
%   \tag{$*$}
% \end{equation*}
% Verifying the three conditions in the definition is an exercise.

Ketika memotivasi definisi hasil kali dalam, kita mengatakan bahwa
perluasan hasil kali titik seharusnya menghasilkan gagasan panjang yang wajar.
Sekarang kita memiliki definisi yang diperlukan untuk menyatakan hal itu secara
tepat.
Ambil
$\norm{\vec{v}}=\sqrt{\smashdp{\innerprod{\vec{v}}{\vec{v}}}}$.
Kita akan menunjukkan bahwa operasi ini memenuhi tiga syarat dalam definisi
norma.
Syarat pertama telah dibahas bersama hasil kali dalam.
Untuk membuktikan syarat kedua,
$\norm{z\cdot\vec{v}}
=\sqrt{\smashdp{\innerprod{z\vec{v}}{z\vec{v}}}}
=\sqrt{\smashdp{z\compconj{z}\cdot\innerprod{\vec{v}}{\vec{v}}}}
=\sqrt{z\compconj{z}}\cdot\sqrt{\smashdp{\innerprod{\vec{v}}{\vec{v}}}}
=\absval{z}\cdot\sqrt{\smashdp{\innerprod{\vec{v}}{\vec{v}}}}$.

Untuk syarat ketiga, seperti dalam bab pertama, kita kembali mengaitkannya
dengan Ketaksamaan Cauchy--Schwarz,\index{Ketaksamaan Cauchy--Schwarz}
$\absval{\innerprod{\vec{v}}{\vec{w}}}\leq \norm{\vec{v}}\cdot\norm{\vec{w}}$.
Pembuktian ketaksamaan di sini memperluas pembuktian pada bab pertama,
tetapi lebih rumit.
Perhatikan selisih $\vec{v}-z\cdot\vec{w}$ ketika $z\in\C$ berubah-ubah.
Norma kuadrat vektor ini merupakan bilangan real tak negatif,
$0\leq \norm{\vec{v}-z\cdot\vec{w}}^2
   =\innerprod{\vec{v}-z\vec{w}}{\vec{v}-z\vec{w}}$.
Linearitas memberikan
$\innerprod{\vec{v}}{\vec{v}-z\vec{w}}
  -z\cdot\innerprod{\vec{w}}{\vec{v}-z\vec{w}}$
dan kemudian
$\innerprod{\vec{v}}{\vec{v}}
  -\compconj{z}\cdot\innerprod{\vec{v}}{\vec{w}}
  -z\cdot\innerprod{\vec{w}}{\vec{v}}
  +z\compconj{z}\cdot\innerprod{\vec{w}}{\vec{w}}$.
Pernyataan ini berlaku untuk setiap $z\in\C$.
\begin{equation*}
 0
  \leq\innerprod{\vec{v}}{\vec{v}}
  -\compconj{z}\cdot\innerprod{\vec{v}}{\vec{w}}
  -z\cdot\compconj{\innerprod{\vec{v}}{\vec{w}}}
  +z\compconj{z}\cdot\innerprod{\vec{w}}{\vec{w}}
\end{equation*}

Ambil
$z=\innerprod{\vec{v}}{\vec{w}}/\innerprod{\vec{w}}{\vec{w}}$
(jika $\vec{w}=\zero$, Ketaksamaan Cauchy--Schwarz berlaku langsung,
sehingga kasus ini kita abaikan).
Penyebutnya merupakan bilangan real, sehingga
$\compconj{z}=\compconj{\innerprod{\vec{v}}{\vec{w}}}/\innerprod{\vec{w}}{\vec{w}}$,
dan substitusinya ke suku kedua memberikan
$\compconj{\innerprod{\vec{v}}{\vec{w}}}\innerprod{\vec{v}}{\vec{w}}/\innerprod{\vec{w}}{\vec{w}}$.
Substitusi ke suku ketiga memberikan hasil yang sama, demikian pula
substitusi ke suku terakhir.
Menggabungkan ketiga suku sejenis itu menghasilkan
\begin{equation*}
 0
  \leq\innerprod{\vec{v}}{\vec{v}}
  -\frac{\innerprod{\vec{v}}{\vec{w}}\compconj{\innerprod{\vec{v}}{\vec{w}}}}{\innerprod{\vec{w}}{\vec{w}}}
  =\norm{\vec{v}}^2
  -\frac{\absval{\innerprod{\vec{v}}{\vec{w}}}^2}{\norm{\vec{w}}^2}
\end{equation*}
Menata ulang dan menghilangkan penyebut memberikan ketaksamaan yang diinginkan.

Dengan hasil itu, kita dapat membuktikan Ketaksamaan Segitiga.
Sekali lagi, pembuktian ini memperluas argumen dalam bab pertama.
\begin{align*}
  \norm{\vec{v}+\vec{w}}^2
  =\innerprod{\vec{v}+\vec{w}}{\vec{v}+\vec{w}}  
  &=\innerprod{\vec{v}}{\vec{v}}
    +\innerprod{\vec{v}}{\vec{w}}  
    +\innerprod{\vec{w}}{\vec{v}}  
    +\innerprod{\vec{w}}{\vec{w}}   \\ 
  &=\norm{\vec{v}}^2
    +\innerprod{\vec{v}}{\vec{w}}  
    +\compconj{\innerprod{\vec{v}}{\vec{w}}}  
    +\norm{\vec{w}}^2              \\   
  &=\norm{\vec{v}}^2
    +2\cdot\RE{\innerprod{\vec{v}}{\vec{w}}}  
    +\norm{\vec{w}}^2      \\
  &\leq\norm{\vec{v}}^2
    +2\absval{\innerprod{\vec{v}}{\vec{w}}}  
    +\norm{\vec{w}}^2     \\
% \end{align*}
\intertext{Terapkan Ketaksamaan Cauchy--Schwarz.}
  &\leq\norm{\vec{v}}^2
    +2\cdot\norm{\vec{v}}\cdot\norm{\vec{w}}
    +\norm{\vec{w}}^2      
  =(\norm{\vec{v}}+\norm{\vec{w}})^2         
\end{align*}
Mengambil akar kuadrat kedua ruas memberikan hasil yang diinginkan.

Mari berhenti sejenak untuk melihat satu lagi hubungan antara hasil kali dalam,
norma, dan geometri.
Berikut diagram jajaran genjang yang sudah dikenal untuk jumlah dua vektor;
selain $\vec{v}+\vec{w}$, kita juga menggambarkan $\vec{v}-\vec{w}$.
\begin{center}
  \includegraphics{jc/asy/innerproduct001.pdf}
\end{center}
Kita akan membuktikan
\definend{identitas jajaran genjang},\index{identitas jajaran genjang}
$\norm{\vec{v}+\vec{w}}^2+\norm{\vec{v}-\vec{w}}^2=2\norm{\vec{v}}^2+2\norm{\vec{w}}^2$.
(Identitas ini memberikan pernyataan yang menarik: pada suatu jajaran genjang,
jumlah kuadrat panjang kedua diagonal sama dengan jumlah kuadrat panjang
keempat sisinya.)
Walaupun gambar di atas memperlihatkan bidang, tempat kita biasanya memakai
notasi hasil kali titik dan menuliskan panjang sebagai $\absval{\vec{v}\,}$,
di sini kita akan memakai notasi hasil kali dalam dan norma.

Seperti di atas, uraikan
$
  \norm{\vec{v}+\vec{w}}^2=\innerprod{\vec{v}+\vec{w}}{\vec{v}+\vec{w}}
$
dengan menggunakan sifat aditif pada masukan pertama, lalu pada masukan kedua.
\begin{align*}
  \norm{\vec{v}+\vec{w}}^2
  &=\innerprod{\vec{v}+\vec{w}}{\vec{v}+\vec{w}}           \\
   % &=\innerprod{\vec{v}}{\vec{v}+\vec{w}}+\innerprod{\vec{w}}{\vec{v}+\vec{w}} \\
  &=\innerprod{\vec{v}}{\vec{v}}            
   +\innerprod{\vec{v}}{\vec{w}}
   +\innerprod{\vec{w}}{\vec{v}}
   +\innerprod{\vec{w}}{\vec{w}}
  =\norm{\vec{v}}^2            
   +\innerprod{\vec{v}}{\vec{w}}
   +\innerprod{\vec{w}}{\vec{v}}
   +\norm{\vec{w}}^2
\end{align*}
Penguraian yang sama untuk $\innerprod{\vec{v}-\vec{w}}{\vec{v}-\vec{w}}$
memberikan
\begin{equation*}
  \norm{\vec{v}-\vec{w}}^2
  =\innerprod{\vec{v}-\vec{w}}{\vec{v}-\vec{w}}
  =\norm{\vec{v}}^2
   -\innerprod{\vec{v}}{\vec{w}}
   -\innerprod{\vec{w}}{\vec{v}}
   +\norm{\vec{w}}^2
\end{equation*}
Menjumlahkan keduanya memberikan hasil yang diinginkan.

Sekarang, dengan operasi hasil kali dalam dan norma, kita dapat menafsirkan
gagasan-gagasan geometri bahkan dalam ruang yang tidak memiliki geometri yang
jelas.
Sebagai contoh, kita dapat memperumum jarak menjadi operasi
\definend{metrik}\index{metrik}
$d(\vec{v},\vec{w})=\norm{\vec{v}-\vec{w}}$.
Contoh lain ialah bahwa kita dapat menyebut vektor $\vec{v}$ dan~$\vec{w}$
ortogonal apabila $\innerprod{\vec{v}}{\vec{w}}=0$, dan kita dapat memperluas
proses Gram--Schmidt ke sebarang ruang hasil kali dalam berdimensi hingga.

Sebuah catatan teknis:~kita telah mendefinisikan operasi norma dari hasil kali
dalam sebagai $\norm{\vec{v}}=\sqrt{\smashdp{\innerprod{\vec{v}}{\vec{v}}}}$,
lalu menurunkan identitas jajaran genjang.
Pada ruang vektor real, sebaliknya juga benar bahwa jika suatu norma memenuhi
identitas ini, maka norma tersebut berasal dari suatu hasil kali dalam, yakni
$\innerprod{\vec{v}}{\vec{w}}
  =(1/2)\cdot(\norm{\vec{v}+\vec{w}}^2-\norm{\vec{v}}^2-\norm{\vec{w}}^2)$.
Namun, pembuktiannya berada di luar cakupan kita.

Suatu norma mungkin tidak berasal dari hasil kali dalam; kita telah melihat
contohnya di atas.
Pada ruang vektor real $\matspace_{\nbyn{2}}$, definisikan norma matriks
$\norm{M}$ sebagai nilai mutlak terbesar dari entrinya.
Perhatikan matriks-matriks berikut.
\begin{equation*}
  M_1=
  \begin{mat}
  1  &2 \\
  3  &4  
  \end{mat}
  \qquad
  M_2=
  \begin{mat}
  8  &7  \\
  6  &5
  \end{mat}
\end{equation*}
Jelas bahwa $\norm{M_1}=4$ dan~$\norm{M_2}=8$.
Mudah pula diperiksa bahwa $\norm{M_1+M_2}=9$ dan
$\norm{M_1-M_2}=7$.
Jadi $\norm{M_1+M_2}^2+\norm{M_1-M_2}^2=130$, sedangkan
$2(\norm{M_1}^2+\norm{M_2}^2)=160$, dan keduanya tidak sama.

Kita menutup bagian ini dengan sebuah catatan tentang definisi hasil kali dalam.
Sebagian penulis, khususnya dalam fisika, mendefinisikan hasil kali dalam agar
linear pada argumen kedua, bukan pada argumen pertama.
Untuk sesuatu yang kita tulis sebagai $\innerprod{\vec{v}}{\vec{w}}$,
para penulis itu menulis $\langle w\, |\, v\rangle$.
Inilah \definend{notasi bra-ket}\index{notasi bra-ket}\index{hasil kali dalam!dibandingkan dengan bra-ket}
yang lazim dalam mekanika kuantum.

\begin{exercises}
\item 
Buktikan bahwa $\R^n$ yang dilengkapi hasil kali titik merupakan ruang hasil
kali dalam.
\begin{answer}
Paragraf pembuka Topik ini menyatakan bahwa hasil kali titik linear pada kedua
masukannya, sehingga syarat pertama terpenuhi.
Syarat~(2) berkaitan dengan konjugasi, yang tidak mengubah bilangan real,
sehingga syarat ini terpenuhi dengan sendirinya.
Syarat~(3) mengharuskan panjang (panjang Euklides biasa pada ruang real
berdimensi~$n$) bernilai nol jika dan hanya jika vektornya adalah~$\zero$.
Jika vektor tersebut mempunyai komponen $v_1$, \ldots~$v_n$, maka
$v_1^2+\cdots+v_n^2=0$ jika dan hanya jika semua komponennya nol.
\end{answer}

\item
Tunjukkan bahwa hasil kali dalam Hermitian pada ruang vektor berdimensi dua
$\C^2$ memenuhi ketiga syarat dalam definisi hasil kali dalam.
\begin{answer}
Pembahasan dalam tubuh Topik membuktikan~(1) dan~(2).
Untuk~(3), misalkan dua anggota $\vec{v},\vec{w}$ dari~$\C^n$
mempunyai komponen ke-$i$ berturut-turut $v_i$ dan $w_i$, sehingga
$\innerprod{\vec{v}}{\vec{w}}=v_1\compconj{w}_1+\cdots+v_n\compconj{w}_n$.
Maka $\innerprod{\vec{v}}{\vec{v}}=v_1\compconj{v}_1+\cdots+v_n\compconj{v}_n$
sama dengan $\absval{v_1}^2+\cdots+\absval{v_n}^2$,
yang jelas bernilai~$0$ jika dan hanya jika setiap komponennya bernilai~$0$.
\end{answer}

\item Tunjukkan bahwa operasi berikut merupakan hasil kali dalam pada ruang
vektor real $\matspace_{\nbyn{2}}$.
\begin{equation*}
  \innerprod{
    \begin{mat}
    a  &b  \\
    c  &d
    \end{mat}
    }{
      \begin{mat}
       e  &f  \\
       g  &h
      \end{mat}
    }
    =
    ae+bf+cg+dh
\end{equation*}
\begin{answer}
Syarat~(1), yakni linearitas pada argumen pertama, mudah dibuktikan.
\begin{multline*}
  \innerprod{
    \begin{mat}
    a_1  &b_1  \\
    c_1  &d_1
    \end{mat}
    +
    \begin{mat}
    a_2  &b_2  \\
    c_2  &d_2
    \end{mat}
    }{
      \begin{mat}
       e  &f  \\
       g  &h
      \end{mat}
    }                                                  \\
  \begin{split}
    &=(a_1+a_2)e+(b_1+b_2)f+(c_1+c_2)g+(d_1+d_2)h   \\
    &=(a_1e+b_1f+c_1g+d_1h)+(a_2e+b_2f+c_2g+d_2h)   \\
    &=
  \innerprod{
    \begin{mat}
    a_1  &b_1  \\
    c_1  &d_1
    \end{mat}
    }{
      \begin{mat}
       e  &f  \\
       g  &h
      \end{mat}
    }
  +
  \innerprod{
    \begin{mat}
    a_2  &b_2  \\
    c_2  &d_2
    \end{mat}
    }{
      \begin{mat}
       e  &f  \\
       g  &h
      \end{mat}
    }
  \end{split}
\end{multline*}
Perkalian skalar juga dapat diperiksa secara langsung.
Syarat~(2) berkaitan dengan konjugasi, yang tidak mengubah bilangan real,
sehingga syarat ini terpenuhi dengan sendirinya.
Untuk~(3), tentu saja $a^2+b^2+c^2+d^2\geq 0$, dengan kesamaan jika dan hanya
jika setiap entrinya nol.
\end{answer}
%
\item Perhatikan ruang vektor real polinomial kuadrat~$\polyspace_2$.
Tunjukkan bahwa operasi berikut merupakan hasil kali dalam:
$\innerprod{a_2x^2+a_1x+a_0}{b_2x^2+b_1x+b_0}=a_2b_2+a_1b_1+a_0b_0$. 
Perhatikan norma yang dihasilkannya.
Berapakah $\norm{x^2-1}$?
\begin{answer}
Syarat~(1) menyatakan bahwa operasi tersebut linear pada masukan pertama.
Pemeriksaan bahwa operasi ini mempertahankan penjumlahan dapat dilakukan
secara langsung:
\begin{multline*}
  \innerprod{(a_2x^2+a_1x+a_0)+(c_2x^2+c_1x+c_0)}{b_2x^2+b_1x+b_0}  \\
  \begin{split}
    &=(a_2+c_2)b_2+(a_1+c_1)b_1+(a_0+c_0)b_0             \\
    &=(a_2b_2+a_1b_1+a_0b_0)+(c_2b_2+c_1b_1+c_0b_0)             \\
    &=\innerprod{(a_2x^2+a_1x+a_0)}{b_2x^2+b_1x+b_0}      \\
      &\qquad+\innerprod{(c_2x^2+c_1x+c_0)}{b_2x^2+b_1x+b_0}
  \end{split}
\end{multline*}
Pemeriksaan bahwa operasi ini mempertahankan perkalian skalar juga dapat
dilakukan secara langsung.
Karena ini ruang vektor real dan konjugasi tidak mengubah skalar bilangan real,
operasi tersebut memenuhi syarat~(2).
Untuk syarat~(3), ambil $\vec{v}=a_2x^2+a_1x+a_0$.
Maka $\innerprod{\vec{v}}{\vec{v}}=a_2^2+a_1^2+a_0^2$ jelas tak negatif dan
sama dengan~$0$ jika dan hanya jika $\vec{v}$ merupakan polinomial nol.

Untuk norma $x^2-1$, kita memperoleh
\begin{equation*} 
 \norm{x^2-1}=\sqrt{\innerprod{x^2-1}{x^2-1}}=\sqrt{1^2+0^2+(-1)^2}=\sqrt{2}
\end{equation*}
\end{answer}

\item Buktikan bahwa himpunan fungsi bernilai real yang kontinu pada interval
tertutup $\closedinterval{0}{1}$ merupakan ruang vektor.
\begin{answer}
Kesepuluh syarat dalam definisi ruang vektor merupakan akibat dari
hasil-hasil dalam Kalkulus.
Sebagai contoh, syarat pertama ialah ketertutupan terhadap penjumlahan, dan
dari Kalkulus kita mengetahui bahwa jumlah dua fungsi kontinu merupakan fungsi
kontinu.
(Meskipun demikian, mahasiswa sering baru melihat pembuktian yang sepenuhnya
teliti dalam mata kuliah berikutnya.)
Demikian pula, penjumlahan dua fungsi bersifat komutatif.
Tentu saja, vektor nolnya adalah fungsi $f(x)=0$, yang kontinu.
\end{answer}

%
\item
Tunjukkan bahwa upaya mendefinisikan $\innerprod{\vec{v}}{\vec{w}}$ pada
$\C^n$ agar linear pada kedua masukan menyebabkan
$\innerprod{\vec{v}}{\vec{v}}$ tidak lagi memberikan panjang.
\begin{answer}
Andaikan $\innerprod{\vec{v}}{\vec{v}}$ lebih besar daripada atau sama
dengan~$0$ untuk setiap vektor kompleks~$\vec{v}$.
Linearitas pada kedua masukan memberikan
$0\leq \innerprod{i\vec{v}}{i\vec{v}}
  =i\innerprod{\vec{v}}{i\vec{v}}
  =i^2\innerprod{\vec{v}}{\vec{v}}\leq 0$.
Jadi $\innerprod{\vec{v}}{\vec{v}}=0$ untuk setiap~$\vec{v}$, yang bertentangan
dengan syarat positif-tegas ketika $\vec{v}\neq\zero$.
\end{answer}


% \item Show that this operation
% on the vector space of continuous real-valued functions 
% over $\closedinterval{0}{1}$ satisfies the 
% conditions in the definition of a norm.  
% \begin{equation*}
%   \norm{f}
%   =\int_{x=0}^1 \absval{f(x)}\,dx
% \end{equation*}
% \begin{answer}
% The first condition holds because the integral of an absolute
% value cannot be negative, and for continuous functions it is zero
% if and only if the function is zero.
% The second condition is  result from elementary Calculus.
% \end{answer}



\end{exercises}
\index{hasil kali dalam|)}
\endinput
% \end{document}
